\section{Proofs for R\'enyi, Tsallis and mixed-state entropies}\label{sec:fqe-spectral-proofs}

\begin{proof}[Proof of Corollary~\ref{cor:fqe-tsallis}]The limit \(q\to1\) is the von Neumann entropy bound \emph{in nats} (divide by
\(\ln2\) for bits). For the proof take \(f_q(x)=(x^q-x)/(1-q)\), with \(f_q(0)=0\). On \((0,1)\), \[
f_q'''(x)=q(2-q)x^{q-3}>0,\qquad
f_q^{(4)}(x)=-q(2-q)(3-q)x^{q-4}<0.
\] Thus its Hermite cubic has \(a_q>0\) and majorises \(f_q\), including \(0\) by continuity. Set
\(g(q)=(15-7q)3^{q-1}/(5+3q)\). Direct divided differences give \[
b_q+\tfrac{17}{12}a_q
=\frac{6^{3-q}}{48(1-q)}
\bigl[5+3q-(15-7q)3^{q-1}\bigr]
=\frac{6^{3-q}(5+3q)}{48}\cdot\frac{1-g(q)}{1-q}<0. \tag{T}\label{eq:fqe-tsallis-coefficient}
\] The prefactor is positive, so the sign claim is that \(g>1\) on \((0,1)\) and \(g<1\) on \((1,2)\).
Indeed \(g(0)=g(1)=1\) and \[
(\log g)'=\ln3-\frac{80}{(15-7q)(5+3q)}.
\] Since \((15-7q)(5+3q)=75+10q-21q^2\), and \(16/15<\ln3<5/4\), this derivative is positive at \(q=0\),
increases up to the vertex \(q=5/21\) of the quadratic and then decreases, so it crosses zero exactly once in
\((0,1)\) (it is \(\ln3-\tfrac54<0\) at \(q=1\)) and remains negative on \([1,2)\). Thus \(g>1\) on \((0,1)\)
and \(g<1\) on \((1,2)\), proving the strict sign in \eqref{eq:fqe-tsallis-coefficient} on either
side of \(q=1\). The logarithm bounds follow, for example, from \(\ln3=\sum_{k\ge0}1/((2k+1)4^k)\):
its first two terms give the lower bound \(13/12\), and using \(2k+1\ge3\) in the terms with \(k\ge2\) gives the upper
bound \(13/12+\tfrac13\sum_{k\ge2}4^{-k}=10/9\). At \(q=1\), since \(g(1)=1\), the quotient \((1-g(q))/(1-q)\) tends to \(g'(1)=(\log g)'(1)=\ln3-\tfrac54\),
so the continuous limit is \(b_1+17a_1/12=\tfrac{36\cdot8}{48}(\ln3-\tfrac54)=6\ln3-\tfrac{15}2<0\).
As for the von Neumann entropy, \(f_q^{(4)}<0\) makes the Hermite remainder vanish on \([0,1]\) only at the two nodes, so
equality can also be identified without the zeros of \(F\): \(P=1\) and \(\Delta_f=0\) give every cut the spectrum
\((\tfrac12,\tfrac16,\tfrac16,\tfrac16)\), and Proposition~\ref{prop:fqe-onecut} applies.

\end{proof}

\paragraph{Sharpness of the Tsallis range.} At \(q=2\) the bound is still valid, but equality is not
unique: \(S_2(\rho)=1-\operatorname{tr}\rho^2\) and Lemma~\ref{lem:fqe-purity} give \(\mathcal
E_{S_2}(\psi)=1-P/3\le2/3\), with equality iff all one-qubit marginals are maximally mixed and
\(\tau=0\). The cluster state $C_4$ of the introduction has cut spectra \((1/2,1/2,0,0)\), \((1/4)^4\), \((1/4)^4\), and so also attains \(2/3\) at
\(q=2\). More strongly, \(M_4\) is not a mean Tsallis-\(q\) maximiser for any \(q>2\). The
averaged cut moments are \[
m_q=2^{-q}+3\,6^{-q},\qquad
c_q=(2^{1-q}+2\,4^{1-q})/3,\qquad
m_q-c_q=\frac{2^{-q}}3(1+9\,3^{-q}-8\,2^{-q}).
\] The last parenthesis vanishes at \(q=2\) and has strictly positive derivative
\(8(\ln2)2^{-q}-9(\ln3)3^{-q}\) for \(q\ge2\): indeed \((3/2)^q\ge9/4>9\ln3/(8\ln2)\) follows from
\(4>3\). Therefore \(S_q(C_4)>S_q(M_4)\) for all \(q>2\). This explains why the parameter range
\(0<q<2\) for unique \(M_4\) maximisers ends sharply at 2.

\paragraph{Expanded spectral gap.} In the setting of Proposition~\ref{prop:fqe-majorant}, with $\alpha$ and $\Delta_f$ as in its proof, the purity identity
(Lemma~\ref{lem:fqe-purity}) and the remainder identity~\eqref{eq:fqe-remainder} at \(N=1\), \(z=\psi\), turn
\eqref{eq:fqe-spectral} into the equivalent formula
\[
\begin{aligned}
3\bigl(\mathcal E_f(M_4)-\mathcal E_f(\psi)\bigr)
={}&\left(\frac\alpha2+\frac{a}{16}\right)W_1+\frac\alpha2\tau+\frac a{32}\sum_i\|V_i-8U_i\|^2\\
&+\frac a8\sum_{i<j}\bigl\|(R_i^{(i)}-R_j^{(j)})z\bigr\|^2+\Delta_f,
\end{aligned}\tag{SF+}\label{eq:fqe-spectral-plus}
\]
in which every term is nonnegative. The residual and pair terms sum to $\frac a{12}\bigl(F-\frac34W_1\bigr)$, which
Theorem~\ref{thm:fqe-general-envelope} bounds below by $\frac a{12}(\tau-\tau^{3/2})$.

\paragraph{Tsallis stability.} For \(0<q<2\), with \(\delta_q=\mathcal E_{S_q}(M_4)-\mathcal
E_{S_q}(\psi)\), the identity \eqref{eq:fqe-spectral-plus}, with \(\alpha_q=-b_q-17a_q/12>0\),
and Theorem~\ref{thm:fqe-general-envelope} give
\[
3\delta_q\ge
\left(\frac{\alpha_q}{2}+\frac{a_q}{16}\right)
W_1+\frac{\alpha_q}{2}\tau+\frac{a_q}{12}\bigl(\tau-\tau^{3/2}\bigr).
\] Here \(\alpha_q,a_q>0\) for every fixed \(q\) in this interval (with the von Neumann limit at \(q=1\)), and
\(\tau-\tau^{3/2}\ge0\). At \(q=2\) the bound becomes the identity \(\delta_2=(W_1+\tau)/6\).

\paragraph{Quantitative consequence for von Neumann entropy.} With \(\delta=E(M_4)-E(\psi)\) in bits,
Theorem~\ref{thm:fqe-general-envelope} gives \[
3\ln2\,\delta\ge
\frac{21-15\ln3}{8}W_1
+\Bigl(\frac{15}4-3\ln3\Bigr)\tau+\frac32(\ln3-1)\bigl(\tau-\tau^{3/2}\bigr).
\] Both \(\tau\) terms are nonnegative, and for small \(\tau\) the
coefficient of \(\tau\) rises from \(\frac{15}4-3\ln3\approx0.454\) to \(\frac94-\frac32\ln3\approx0.602\). Dropping $F$ in \eqref{eq:fqe-entropy-gap} gives the simpler joint estimate \[
W_1+\tau
\le\frac{6\ln2}{\tfrac{15}{2}-6\ln3}\,\delta
\approx4.57862\,\delta.
\] Thus near-maximal entropy forces small one-qubit polarisation and spin-flip overlap.
These conditions alone do not identify the state: cluster states also have $W_1=\tau=0$.
Appendix~\ref{sec:fqe-orbit-stability} adds control of the distance to the two HS orbits and transfers it to
other strict spectral majorants.

\begin{proof}[Proof of Corollary~\ref{cor:fqe-convex-roof}]
The ceiling is proved in the main text after the corollary; it remains to show that no mixed state attains it. Fix one \(2|2\) cut. The support of a
mixed \(\rho\) contains independent vectors with \(4\times4\) flattenings
\(A,B\). Their pencil contains a nonzero singular matrix: either a
generator is singular, or both are invertible and the positive-degree
polynomial \(\det(A+tB)\) has a complex root. Let \(\chi\) be the
normalised support vector represented by that matrix. Because \(\rho\)
is positive definite on its support, for some \(p>0\) the operator
\(\rho-p|\chi\rangle\langle\chi|\) is positive. Thus an ensemble for
\(\rho\) can contain \(\chi\) with positive weight. But \(\chi\) has
deficient Schmidt rank, whereas every pure maximiser has full rank on
this cut. Applying the strict pure-state bound to that ensemble makes
the mean of the convex roofs strictly smaller than the ceiling.

For R\'enyi entropies in the stated range, use the pure-state bound of
Corollary~\ref{cor:fqe-renyi-gw} in place of
\eqref{eq:fqe-spectral}; the same support argument proves strictness.
\end{proof}

The same exclusion of mixed maximisers works whenever all pure maximisers have full Schmidt rank on
one fixed square bipartition.
